\documentclass[11pt,a4paper]{article}
\usepackage[margin=24mm]{geometry}
\usepackage{amsmath,amssymb,amsthm,hyperref}
\newtheorem{theorem}{Theorem}\newtheorem{lemma}[theorem]{Lemma}\newtheorem{proposition}[theorem]{Proposition}
\newcommand{\gen}[1]{\langle #1\rangle}
\title{Coset trialities, chamber recognition and reduced relations}
\author{Human Compression project: recovery of L122}\date{}
\begin{document}\maketitle
Leemans and Stokes~\cite{LS} construct coset geometries with trialities and
compare their automorphism groups with those of reduced incidence graphs.
We give the complete elementary interfaces needed by their finite examples.
The subgroup construction and reduction are their ideas; the order and graph
conventions are made explicit. A type-permuting correlation need not have
order equal to the length of its type cycle.

\section{The subgroup action and its exact order}
Let $G$ be a group and $(H_0,H_1,H_2)$ an indexed subgroup family.
The objects of type $i$ are the right cosets $H_i g$; objects of different
types are incident when their underlying cosets intersect. The right action
of $G$ preserves incidence. No assertion that it is the full automorphism
group follows merely from this construction.

\begin{proposition}\label{cycle}
Suppose $\sigma\in\operatorname{Aut}G$ maps $H_i$ to $H_{i+1}$, with indices
modulo three. Then
\[
 \tau(H_i g)=H_{i+1}\sigma(g)
\]
is an incidence-preserving type cycle. It has order three if and only if
$H_i\sigma^3(g)=H_i g$ for every $i,g$. In particular, $\sigma^3=1$ suffices.
\end{proposition}
\begin{proof}
The map is well defined because it sends every equality of cosets through a
group automorphism. It is bijective, and a common coset element maps to a
common coset element; applying the inverse proves the converse incidence
implication. Its cube sends $H_i g$ to $H_i\sigma^3(g)$, giving the exact
criterion. The type cycle rules out orders one and two.
\end{proof}
The same proof works for any number $r$ of cyclically permuted types, with
$\sigma^r$ in place of $\sigma^3$. The requirement that the pregeometry is
a geometry remains separate. Lemma~1 and Method~1 of~\cite{LS} need this
order qualification to identify the \emph{induced} correlation as a triality
under their initial definition. We do not infer that their stronger bare
existence assertion is false solely from the omitted order argument.

\section{A complete rank-three recognition criterion}
\begin{theorem}\label{recognition}
Let $G=\gen{r_0,r_1,r_2}$, where the $r_i$ are distinct involutions, and put
$H_i=\gen{r_j,r_k}$ for $\{i,j,k\}=\{0,1,2\}$. Suppose
\[
 H_i\cap H_j=\gen{r_k},\qquad \bigcap_iH_i=1,
 \qquad H_0\cap H_1H_2=(H_0\cap H_1)(H_0\cap H_2). \tag{1}
\]
Then the coset system is a thin, residually connected, flag-transitive
geometry. Its full type-preserving automorphism group is the given right
copy of $G$.
\end{theorem}
\begin{proof}
Every incident pair of cosets has a common element, and so is a translate
of the corresponding pair of standard parabolics. Consider a triple of
pairwise incident cosets. Translate the intersection of its last two cosets
to the identity, leaving the triple $H_0g,H_1,H_2$. Choose
$h_1\in H_0g\cap H_1$ and $h_2\in H_0g\cap H_2$.
Then $h_1h_2^{-1}\in H_0\cap H_1H_2$ and (1) writes it as $uv$ with
$u\in H_0\cap H_1$, $v\in H_0\cap H_2$.
The element $u^{-1}h_1=vh_2$ belongs to all three cosets. Thus every
chamber has a common element and is a translate of $(H_0,H_1,H_2)$.
Every incident pair extends to such a chamber, as does every singleton.
This proves the geometry and flag-transitivity assertions.

The chamber stabilizer is $\bigcap H_i=1$, so $G$ acts regularly on chambers.
A panel missing type $i$ has exactly $|H_j\cap H_k|=2$ completions.
The rank-two residue at $H_i$ is the coset graph of $H_i$ with the two
subgroups $\gen{r_j}$ and $\gen{r_k}$, which generate $H_i$; it is connected.
This last connectivity can be seen directly by writing a group element as
an alternating product of the two subgroup generators and following the
corresponding incident cosets. The full incidence graph is connected for
the same reason, since the three $H_i$ generate $G$. These are all residues
of rank at least two.

The two chambers in the panel of the standard chamber missing type $i$
are the standard chamber and its translate by $r_i$. The chamber adjacency
graph is therefore the connected Cayley graph generated by the $r_i$.
A type-preserving automorphism fixing one chamber fixes its unique neighbor
of each type, and induction along this graph fixes every chamber and object.
Any type-preserving automorphism can first be composed with an element of
$G$ to fix the standard chamber. It is then the identity. This proves the
full automorphism assertion, without a graph-automorphism census.
\end{proof}
In the finite certificates, each side on the right of (1) has order two
and they intersect trivially. Their product has four elements and is always
contained in the left side. Thus computing that the left side has four
elements proves equality. This is a proof premise, not a heuristic order
comparison. The certificates check all three cyclic versions as well as all
64 subgroup intersections of the marked triple.

\section{Lossless reduction for an actual triality}
Assume that $\tau$ is an incidence-preserving correlation of order three
cycling all types of a rank-three pregeometry. Let $V$ be its type-zero
objects and define the directed relation
\[
 A(x,y)\quad\Longleftrightarrow\quad x\text{ is incident with }\tau(y).
 \tag{2}
\]
Loops are permitted. This is the reduced oriented graph, with the absolute
vertices marked by its loops. In the graphical convention of~\cite{LS},
these marks are called half-edges.

\begin{theorem}\label{reduced}
Relation (2) reconstructs the entire typed pregeometry. Its automorphism
group is naturally the centralizer of $\tau$ in the full type-preserving
automorphism group of the pregeometry.
\end{theorem}
\begin{proof}
Every object has a unique expression $\tau^i(x)$ with $x\in V$ and
$i\in\{0,1,2\}$. For different types $i,j$, their incidence is $A(x,y)$
when $j-i=1$ modulo three, and is $A(y,x)$ when $j-i=2$.
Indeed apply $\tau^{-i}$; in the second case use symmetry and
$\tau^3=1$. This proves reconstruction.
An automorphism $f$ of (2) extends by
$\widehat f(\tau^i x)=\tau^i f(x)$. The incidence formulas prove that this
extension is a type-preserving automorphism commuting with $\tau$.
Conversely such an automorphism restricts to an automorphism of (2), and
commutation makes its action everywhere determined by that restriction.
These operations are inverse group homomorphisms.
\end{proof}
For the coset action of Proposition~\ref{cycle}, conjugation by $\tau$ on the
faithful right copy of $G$ induces $\sigma$. If Theorem~\ref{recognition}
applies, Theorem~\ref{reduced} consequently identifies the full reduced
automorphism group with $G^\sigma$. In particular, computing only that
subgroup would not suffice without the full typed-group argument.

\section{Darts, voltages and higher-order qualifications}
For an order-$r$ correlation cycling $r$ types, use coordinates
$(x,i)=\sigma^i x$. Record a directed incidence from $(x,0)$ to $(y,d)$
by the voltage $d\in\mathbb Z/r\mathbb Z$. Its reverse record is $(y,x,-d)$.
The original incidence graph is reconstructed by lifting it to all edges
$(x,i)$--$(y,i+d)$. Every vertex and every \emph{dart} orbit has size $r$;
vertices and darts therefore divide in number by $r$. This proves the
quotient-cover assertion in the dart convention, including its exact lift.
An undirected edge orbit can instead have size $r/2$ when a half-turn
interchanges its endpoints. Such an orbit is a quotient semiedge. A quotient
loop with nonzero voltage is a different object: for $r=3$ an absolute
vertex lifts to a triangle, so its mark must retain that voltage-loop meaning.
The single-dart semiedge convention appropriate to a polarity is not an
ordinary graph-cover interpretation of this triality mark.

The order hypothesis cannot be removed by redefining the correlation on
its last type, as suggested after Lemma~4 of~\cite{LS}. Here are elementary
counterexamples to that suggested normalization; neither changes the
order-three examples of the application.

\begin{proposition}\label{normalization}
There are rank-two and rank-three geometries with cyclic type-permuting
correlations for which resetting the last type map to make its cube or
square the identity does not preserve incidence.
\end{proposition}
\begin{proof}
For rank two take the cycle on $\mathbb Z/6\mathbb Z$, with the even and odd
vertices as types, and $\delta(x)=x+1$. The proposed reset agrees with
$\delta$ on even vertices and with $\delta^{-1}$ on odd vertices.
It sends the incident pair $(1,2)$ to $(0,3)$, which is not incident.
The reduced directed relation on the even vertices has loops and a directed
three-cycle. Discarding its orientation identifies the two opposite
relations and is not the claimed lossless unlabeled reduction.

For rank three take $\mathbb Z/9\mathbb Z$, types modulo three, and incidence
between vertices whose cyclic distance is one or two. Every edge lies in a
consecutive triple, so this is a geometry. Each vertex residue is a path on
four vertices, so it is residually connected. Rotation $\sigma(x)=x+1$ has
order nine. Resetting the last type map replaces its value on type two by
$\sigma^{-2}$. It again sends the incident pair $(2,3)$ to $(0,4)$,
which is not incident. The natural degree-three voltage lift also fails
without the order condition: the reduced voltage-two record from $3$ to $0$
represents the edge $(3,2)$, whereas its sheet-one lift would be $(4,0)$,
a nonedge. Thus the general cover assertion of Lemma~4 needs the same
qualification, not merely its identification as a quotient.
For completeness its full graph automorphism group is dihedral of order 18:
among its edges those of distance one have two common neighbors, whereas
those of distance two have one; hence its automorphisms preserve the
underlying nine-cycle, and every dihedral symmetry is an automorphism.
Every element of order three is rotation by three or six, preserving each
type. Thus even the existence of an order-three type cycle does not follow
for arbitrary geometries from the existence of an order-nine type cycle.
This last geometry is not being asserted to be the coset system of
Proposition~\ref{cycle}; it is not a counterexample to that lemma's stronger
bare existential formulation.
\end{proof}
For an arbitrary correlation cycling types, one may still retain every
oriented labeled relation between the selected coordinate types. Those
relations reconstruct incidence by applying the actual powers of the
correlation. They do not justify imposing a new order or discarding the
extra relations. The repaired order-$r$ version above is the one used here. More explicitly,
for arbitrary order define $A_d(x,y)$ by $x*\sigma^d(y)$, for
$1\leq d<r$. If $i<j$, incidence of $\sigma^i x$ and $\sigma^j y$
is exactly $A_{j-i}(x,y)$, by applying $\sigma^{-i}$. Thus all the
labeled relations are lossless without any order assumption; what fails is
the proposed order-$r$ voltage lift or the deletion of the extra data.
In rank two a single directed relation suffices. Different dualities can
give different reductions: on the six-cycle, rotation $x\mapsto x+1$
has three reduced loops, whereas the polarity $x\mapsto1-x$ has two.
This also proves the nonuniqueness qualification of source Lemma~2.

\begin{thebibliography}{9}
\bibitem{LS} D. Leemans and K. Stokes, \emph{Coset geometries with trialities
and their reduced incidence graphs}, Acta Math. Univ. Comenianae
\textbf{88}(3) (2019), 911--916.
\bibitem{Artzy} R. Artzy, \emph{Self-dual configurations and their Levi graphs},
Proc. Amer. Math. Soc. \textbf{7} (1956), 299--303.
\bibitem{PS} T. Pisanski and B. Servatius, \emph{Configurations from a Graphical
Viewpoint}, Birkh\"auser, 2013. Quotient and voltage terminology; the incidence
and reconstruction arguments used here are given in full.
\end{thebibliography}
\end{document}
